%------------------------------------------------------------------------------
% Chapter 2: Background and Related Work
%------------------------------------------------------------------------------

\chapter{Background and Related Work}\label{chap:background}

Continuous black-box optimization considers problems in which an algorithm
must make decisions from objective-function evaluations without access to the
analytical structure of the objective. In practical applications, these
evaluations may also be affected by uncertainty, so that the observed fitness
does not always represent the underlying objective exactly. This chapter
introduces the concepts needed to study this setting and reviews previous work
on automated algorithm design using large language models (LLMs).

The chapter first introduces continuous black-box optimization and the BBOB
benchmark suite, together with commonly used derivative-free and evolutionary
optimization methods. It then discusses noisy fitness evaluation, the main
forms of fitness noise, and its effects on optimization algorithms. The
discussion then moves from conventional algorithm design to Automated
Algorithm Design (AAD), program synthesis, and LLM-based evolutionary
approaches. Particular attention is given to LLaMEA, which uses an LLM as a
semantic variation operator within a $(1+1)$ evolutionary search over
executable optimization algorithms. Finally, the effect of noisy fitness on
elitist selection is examined. This provides the basis for the research gap
addressed in this thesis.

%------------------------------------------------------------------------------
\section{Continuous Optimization}\label{sec:continuous_optimization}
%------------------------------------------------------------------------------

Continuous black-box optimization seeks the global minimum of an unknown
real-valued objective function:
\begin{equation}\label{eq:bbob_min}
\min_{x \in \mathcal{S} \subseteq \mathbb{R}^D} f(x)
\end{equation}
where $\mathcal{S}$ denotes a bounded $D$-dimensional continuous search space.
In the BBOB benchmark setting, the search domain is commonly defined within
$[-5,5]^D$. The analytical gradients $\nabla f(x)$ and the internal
formulation of the objective are not available to the optimizer. Instead, the
algorithm has to learn about the search landscape from objective-function
evaluations.

This setting occurs in scientific and engineering problems where an explicit
mathematical model is unavailable, too expensive to evaluate directly, or not
accurate enough for optimization. The No Free Lunch theorem also shows that
no single optimization heuristic can outperform all alternatives on every
possible problem class without suitable problem-specific inductive bias
\cite{Wolpert1997}. The characteristics of the optimization
problem therefore matter when choosing or designing an optimization method.

Several properties can make a continuous black-box problem difficult. These
include separability, conditioning, modality, and interactions between
decision variables. In a separable function, the variables can be optimized
independently, whereas non-separable functions contain dependencies between
variables. Ill-conditioned functions have directions with substantially
different sensitivities, which can make progress difficult for methods that do
not adapt their search geometry. Multimodal functions contain multiple local
optima and therefore require the optimizer to balance exploration of the search
space with exploitation of promising regions.

Because optimization performance depends on these characteristics,
standardized benchmark suites are useful for comparing algorithms under
reproducible conditions. The Black-Box Optimization Benchmark (BBOB) suite
provides a collection of continuous test functions with different landscape
structures \cite{BBOB2009}. The suite is integrated into
experimental frameworks such as \texttt{IOHexperimenter} \cite{deNobel2024}
and COCO. The known optima of the benchmark functions also make it possible to
measure the progress of an optimizer relative to a known reference solution.

The 24 noiseless BBOB functions are organized into five broad groups according
to their structural characteristics. Table~\ref{tab:bbob_groups} summarizes
these groups.

\begin{table}[htbp]
\centering
\caption{Taxonomy of the 24 Continuous Black-Box Optimization Benchmark (BBOB) Functions \cite{BBOB2009}.}
\label{tab:bbob_groups}
\small
\begin{tabularx}{\textwidth}{c c l X}
\toprule
\textbf{Group} & \textbf{Function Indices} & \textbf{Landscape Property} & \textbf{Core Topological Challenge} \\
\midrule
1 & $f_1 - f_5$ & Separable Functions & Orthogonal dimensions; can be optimized independently along coordinate axes. \\
2 & $f_6 - f_9$ & Low/Moderate Conditioning & Functions with low or moderate conditioning, including both unimodal and multimodal structures. \\
3 & $f_{10} - f_{14}$ & High Conditioning & Steep, narrow ill-conditioned valleys requiring appropriate scaling and rotation. \\
4 & $f_{15} - f_{19}$ & Multi-modal with Global Structure & Multiple local optima combined with an identifiable global structure. \\
5 & $f_{20} - f_{24}$ & Multi-modal with Weak Structure & Highly deceptive local optima with weak global structure. \\
\bottomrule
\end{tabularx}
\end{table}

BBOB performance is commonly presented using empirical cumulative
distribution functions (ECDFs). An ECDF shows the proportion of problem
instances that reach a specified target within a given number of function
evaluations \cite{Hansen2010}. A commonly used target in the noiseless
benchmark setting is an optimality gap of
\begin{equation}
\Delta f = f(x)-f(x^*) \leq 10^{-8}.
\end{equation}
Such evaluations consider both the quality of the final solution and the
number of objective evaluations required to reach a given level of precision.

A number of derivative-free and evolutionary methods are commonly used for
continuous black-box optimization. Nelder--Mead maintains a simplex of
candidate points and updates it through operations such as reflection,
expansion, contraction, and shrinkage \cite{Nelder1965}. Since it does not
require derivatives, it can be used when only objective-function evaluations
are available.

Differential Evolution (DE) is a population-based method that generates new
candidate vectors using differences between existing population members
\cite{Storn1997}. Mutation and crossover allow the method to explore the
continuous search space while making use of information contained in the
current population. Particle Swarm Optimization (PSO) uses a population of
particles whose movement depends on their previous positions and information
obtained from the population \cite{Kennedy1995}.

Evolution Strategies form another important class of continuous optimization
methods. They use stochastic variation and selection to update candidate
solutions rather than relying on a fixed geometric search structure. The
Covariance Matrix Adaptation Evolution Strategy (CMA-ES) is particularly
important because it adapts the covariance matrix of its search distribution
according to the observed performance of candidate solutions
\cite{Hansen2001}. This allows the search distribution to adapt to the
geometry of the objective landscape and makes CMA-ES effective on problems
with ill-conditioning and non-separability. It has consequently become a
common reference method for continuous black-box optimization.

NEWUOA provides another derivative-free approach based on local quadratic
models \cite{Powell2006}. Comparisons on difficult continuous landscapes have
shown the advantages of adaptive population-based methods such as CMA-ES on
problems where the local geometry changes substantially.

The design of these algorithms has traditionally relied on human expertise.
One response to this limitation has been the development of modular algorithm
frameworks, in which predefined algorithmic components can be selected and
combined automatically \cite{vanRijn2016,deNobel2021}. Modular CMA-ES
\cite{vanRijn2016} and Modular DE \cite{deNobel2021}, for example, decompose
optimization algorithms into interchangeable components such as covariance
adaptation, sampling mechanisms, and stopping criteria. Automated configurators
can then search through many possible combinations. The search space is,
however, still determined by the components provided by the designer. Such a
system can select and combine existing mechanisms, but introducing an
algorithmic structure outside the predefined component library is more
difficult.

The limitations of deterministic optimization become particularly relevant in
physical systems, where objective values are often obtained from measurements.
Battery-management studies provide examples of optimization and
state-estimation problems based on noisy measured quantities, including
electrical and thermal effects \cite{Plett2004,Xiong2018}. These examples
illustrate how measurement uncertainty can affect model-based decision
processes. A similar situation occurs in finite-element model updating, where
vibration measurements are used to identify or update structural models and
are subject to measurement uncertainty \cite{Mottershead1993,Marwala2010}.
These examples illustrate why an optimization algorithm may not always receive
exactly the same objective value for a given candidate solution.

%------------------------------------------------------------------------------
\section{Noisy Optimization}\label{sec:noisy_optimization}
%------------------------------------------------------------------------------

In practical optimization, repeated evaluations of the same candidate
solution do not necessarily return the same fitness value. Stochasticity can
come from measurement errors, environmental effects, simulation randomness,
numerical effects, or other sources of uncertainty. The optimizer therefore
does not always observe the underlying objective directly.

Jin and Branke classify uncertainty in optimization according to where it
enters the optimization process, including the parameter space, time space,
surrogate-model space, and fitness space \cite{Jin2005}. The focus here is on
fitness-space uncertainty. In this case, the candidate solution itself is
unchanged, while the objective value returned to the optimizer is perturbed.

A general additive representation is
\begin{equation}\label{eq:noisy_evaluation_bg}
f_{\text{noisy}}(x) = f_{\text{clean}}(x) + z,
\end{equation}
where $f_{\text{clean}}(x)$ is the underlying objective and $z$ is a random
disturbance. The observed value $f_{\text{noisy}}(x)$ is therefore a random
variable, even when the same coordinate vector $x$ is evaluated repeatedly.

Different probability distributions can be used to model fitness noise.
Common choices include Uniform, Cauchy, and Gaussian distributions
\cite{Hansen2009Noisy}:

\begin{enumerate}
\item \textbf{Additive Uniform Noise:}
The disturbance is bounded within a fixed interval:
\begin{equation}
z \sim \mathcal{U}(-a,a).
\end{equation}

\item \textbf{Additive Cauchy Noise:}
The Cauchy distribution introduces heavy-tailed deviations:
\begin{equation}
p(z)=
\frac{1}{\pi\gamma\left[1+(z/\gamma)^2\right]}.
\end{equation}

\item \textbf{Additive Gaussian Noise:}
The disturbance is normally distributed:
\begin{equation}
z \sim \mathcal{N}(0,\sigma^2).
\end{equation}
\end{enumerate}

These distributions provide generic examples of fitness-noise
models. The exact noise definitions used in benchmark suites can differ from
these simple additive formulations. Uniform noise limits the disturbance to a
fixed range, whereas Cauchy noise can produce occasional large deviations
because of its heavy tails. Gaussian noise gives a continuous model in which
the variance can be controlled through $\sigma$. The choice of noise model
therefore affects both the frequency and magnitude of deviations between
observed and underlying fitness.

The effect of noise becomes especially important when an algorithm makes
decisions based on relative fitness. Consider two candidate solutions for
which
\begin{equation}
f_{\text{clean}}(x_1) < f_{\text{clean}}(x_2).
\end{equation}
According to the clean objective, $x_1$ is the better solution. A particular
noise realization can nevertheless result in
\begin{equation}
f_{\text{noisy}}(x_1) >
f_{\text{noisy}}(x_2).
\end{equation}
The observed ranking is then reversed. An optimizer that relies on this
ranking may consequently make a selection decision that differs from the one
it would make using the clean objective.

The consequences depend on the optimization method. Direct
search and evolution-strategy methods can respond differently to noisy
evaluations, and noise can make objective differences unreliable when the
search approaches the scale of the perturbations
\cite{ArnoldB03,Arnold2003}. For direct-search methods,
small objective differences can become difficult to distinguish from
evaluation noise, which can affect decisions about the local search
direction or step size. As the search radius decreases, the actual objective
differences can become comparable to the noise magnitude. The resulting
decisions may reduce the search radius even though the current solution is
still far from the optimum.

Population-based Evolution Strategies can provide some robustness by using
information from multiple candidate solutions. For example, intermediate
recombination in multi-parent $(\mu/\mu,\lambda)$-ES can average information
across parent vectors and help maintain productive mutation step sizes under
noisy evaluations \cite{Arnold2003}.

CMA-ES provides a more direct example of the effect of noisy ranking. Its
population update depends on the ordering of candidate solutions according to
their measured fitness \cite{Hansen2001}. The generation mean can be written as
\begin{equation}\label{eq:cma_mean_update_bg}
m^{(g+1)} =
m^{(g)} +
\sum_{i=1}^{\mu}
w_i
\left(
x_{i:\lambda}-m^{(g)}
\right),
\end{equation}
where
\begin{equation}
w_1 \geq w_2 \geq \dots \geq w_\mu > 0
\end{equation}
are the weights associated with candidates ordered according to their
measured fitness,
\begin{equation}
f(x_{1:\lambda}) \leq \dots \leq f(x_{\lambda:\lambda}).
\end{equation}

When the fitness evaluations are noisy, this ranking can be corrupted. A
candidate with a worse clean objective may receive a better noisy evaluation
and consequently a larger selection weight. These rank inversions can affect
both cumulative step-size adaptation and covariance-matrix learning
\cite{Jin2005,Hansen2009Noisy}. Noise can therefore
influence not only which candidate is selected, but also the search
distribution used in later generations.

Repeated evaluations are one common way to reduce this problem. If a
candidate is evaluated $N$ times and the observations are averaged, the
variance of the resulting estimate can be reduced. The additional evaluations
also increase the computational cost, however, creating a trade-off between
more reliable fitness estimates and the available evaluation budget
\cite{Jin2005}.

Gradient-based methods can be affected for a similar reason. Finite-difference
gradient estimates depend on differences between objective evaluations. If
these evaluations contain stochastic errors, the estimated gradient can
deviate substantially from the gradient of the underlying objective. This can
lead to unstable or inefficient update trajectories
\cite{Sarafian2024}.

Overall, fitness noise does not only influence the final solution produced by
an optimizer. It can change the information available to the search process
at every iteration. An optimization algorithm designed and evaluated under
deterministic assumptions may therefore behave differently once its objective
evaluations become stochastic.

%------------------------------------------------------------------------------
\section{Automated Algorithm Design}\label{sec:automated_algorithm_design}
%------------------------------------------------------------------------------

The difficulty of manually designing and configuring optimization algorithms
has motivated research into Automated Algorithm Design (AAD). Early AAD
approaches used Genetic Programming and hyper-heuristics to construct or
select algorithmic procedures automatically \cite{Burke2013}. Instead of
manually specifying every algorithmic decision, these methods search over a
predefined representation of possible solutions.

Modular algorithm design developed this idea further by decomposing existing
optimization algorithms into interchangeable components
\cite{vanRijn2016,deNobel2021}. This allows a large number of algorithm
configurations to be explored without requiring each combination to be
implemented manually. At the same time, the available components and the
representation of the algorithm still constrain what can be generated.

Program synthesis provides a broader search space because the system can
generate executable program structures rather than only selecting from a
fixed set of components. The development of large language models has made
this approach increasingly practical, since LLMs can generate and modify
source code from natural-language instructions.

Romera-Paredes et al.~introduced FunSearch, which combines a language model
with an automated evaluation process to search program space
\cite{Romera2023}. The framework constrains the LLM to generate particular
program components within a larger structure and uses evaluation results to
identify promising programs. This work showed that LLM-generated program
variations can discover algorithmic solutions that perform competitively with
human-designed approaches.

Other work has combined LLM generation with evolutionary search in different
ways. EvoLLM uses language models to generate candidate solutions in numerical
optimization settings \cite{Lange2024}. In contrast to methods that generate
complete optimization algorithms, such approaches can use the model to
propose candidate points or numerical solutions.

Other approaches operate directly on algorithmic code. Evolution of Heuristics
(EoH) uses LLMs to generate and evolve heuristic procedures
\cite{Liu2024}. Algorithm Evolution using LLMs (AEL) combines LLM-based
generation with evolutionary mechanisms to construct heuristic algorithms
\cite{AEL2024}. Reflective Evolution (ReEvo) adds reflection and evolutionary
operations to the generation process \cite{ReEvo2024}. These approaches treat
the LLM as part of an iterative search process rather than using it only as a
one-shot code generator.

The distinction between numerical and programmatic generation is important.
A system that generates numerical coordinate vectors searches directly in the
solution space. A system that generates executable algorithms instead searches
over programs. This program-space search can produce different combinations of
operations and algorithmic structures rather than choosing only among
predefined numerical actions.

Zhang et al.~\cite{ReflectiveSurvey2025} describe several functional roles for
LLMs in evolutionary computation. Selection operators rank candidates using
execution or evaluation metrics. Variation operators modify existing
solutions through semantic mutation or crossover. Reflective operators
provide information from execution, such as compiler errors, tracebacks, and
runtime observations, to the LLM. These roles are useful for understanding
LLM-based AAD because the language model is part of a larger feedback loop
rather than operating independently.

Table~\ref{tab:llm_aad_comparison} summarizes the main LLM-driven AAD
approaches considered in this work.

\begin{table}[htbp]
\centering
\caption{Comparative Taxonomy of LLM-Driven Automated Algorithm Design (AAD) Frameworks.}
\label{tab:llm_aad_comparison}
\small
\setlength{\tabcolsep}{3pt}
\renewcommand{\arraystretch}{1.15}
\begin{tabularx}{\textwidth}{
    >{\raggedright\arraybackslash}p{1.65cm}
    >{\raggedright\arraybackslash}p{2.15cm}
    >{\raggedright\arraybackslash}p{2.15cm}
    >{\raggedright\arraybackslash}p{2.45cm}
    >{\raggedright\arraybackslash}X
}
\toprule
\textbf{Framework} &
\textbf{Output Modality} &
\textbf{Search Space} &
\textbf{Feedback Type} &
\textbf{Reported Evaluation Setting} \\
\midrule

FunSearch \cite{Romera2023}
& Program Skeletons
& Program Space
& Score Metrics
& Deterministic benchmark evaluation \\

EvoLLM \cite{Lange2024}
& Coordinate Vectors
& Continuous
& Point Distributions
& Numerical optimization evaluation \\

EvoPrompt \cite{Guo2023}
& Prompt Strings
& Prompt Space
& Evolutionary Score
& Task-dependent evaluation \\

EoH \cite{Liu2024}
& Code Blocks
& Program Space
& Reflective Prompts
& Reported deterministic evaluation \\

AEL \cite{AEL2024}
& Code Blocks
& Program Space
& Evolutionary Crossover
& Reported deterministic evaluation \\

ReEvo \cite{ReEvo2024}
& Code Blocks
& Program Space
& AST and Crossover
& Reported deterministic evaluation \\

AutoEP \cite{AutoEP2026}
& Hyperparameter Strategies
& Structured Parameter Space
& Zero-Shot ELA Features
& Controlled benchmark evaluation \\

LLaMEA-SAGE \cite{LLaMEASAGE2026}
& Python Classes
& Program Space
& AST and XAI Scaffolding
& Noiseless benchmark evaluation \\

LLaMEA-BO \cite{LLaMEABO2025}
& Python Classes
& Program Space
& Surrogate Scores
& Benchmark evaluation \\

LLaMEA \cite{LLaMEA2024}
& Python Classes
& Program Space
& Tracebacks and Score
& Noiseless benchmark evaluation \\

\textbf{This Thesis}
& \textbf{Python Classes}
& \textbf{Program Space}
& \textbf{Tracebacks and Scaled Score}
& \textbf{Stochastic BBOB evaluation} \\

\bottomrule
\end{tabularx}
\end{table}

The literature also differs in the information provided to the LLM before
generation. AutoEP, for example, extracts Exploratory Landscape Analysis
(ELA) features and uses them to guide hyperparameter configuration and
algorithm control \cite{AutoEP2026}. This gives the model information about the
characteristics of the optimization problem before a candidate configuration
or control strategy is generated.

Other approaches make use of execution feedback. Reflexion shows how runtime
observations and feedback can be incorporated into later model generations
\cite{Shinn2023}. Chen et al.~\cite{Chen2023} found that execution feedback
can improve Python function generation accuracy while reducing the number of
generation samples required. For executable algorithm synthesis, this type of
feedback can include numerical performance as well as syntax errors, runtime
exceptions, and other execution information.

Prompt scaffolding provides another way to constrain the generation process.
EvoPrompt treats prompts themselves as objects of evolutionary optimization
\cite{Guo2023}. Promptomatix introduces structured constraints and schemas for
controlling generation \cite{Murthy2025}. These approaches show that the
prompt can form part of the search mechanism rather than simply being a fixed
instruction. Their direct use as variation mechanisms for continuous
optimization algorithm synthesis under stochastic fitness is not the focus of
the reported evaluations and therefore does not establish their effectiveness
in that setting.

Structural information can also be used when generating source code.
LLaMEA-SAGE extends the LLaMEA approach with Abstract Syntax Tree (AST) and
Explainable AI (XAI) information to characterize generated algorithms
\cite{LLaMEASAGE2026}. Van Stein \cite{vanStein2025} investigates explainable
AAD by relating generated code modifications to characteristics of the
optimization problem. These approaches suggest that algorithm synthesis can
make use of information beyond a single scalar fitness value.

Despite these developments, the evaluation setting is an important common
point across much of this literature. The approaches summarized in
Table~\ref{tab:llm_aad_comparison} primarily report evaluations under
deterministic or otherwise controlled benchmark conditions. Consequently, the
effect of stochastic fitness on the evolutionary selection process used by
these systems remains less explored.

%------------------------------------------------------------------------------
\section{Selection Under Noise}\label{sec:selection_under_noise}
%------------------------------------------------------------------------------

Evolutionary Computation provides the conceptual foundation for many of the
automated search processes discussed above \cite{Eiben2015}. Classical
Evolutionary Algorithms repeatedly perform initialization, variation, fitness
evaluation, and selection. An initial population can be represented as
\begin{equation}
\mathcal{P}_0 = \{\theta_1, \dots, \theta_\mu\},
\end{equation}
where $\mu$ candidate solutions are sampled from the search space. Variation
operators generate new candidates, fitness evaluation assigns a scalar score,
and selection determines which candidates are retained for subsequent
generations \cite{Eiben2015}.

The resulting search trajectory depends on the balance between exploration
and exploitation \cite{Goldberg1989}. Variation introduces new candidate
solutions and allows the search to explore previously unvisited regions,
while selection favors candidates that have performed well and concentrates
future search around them.

Population-based Evolution Strategies can reduce the influence of individual
fitness observations by considering multiple candidates. A $(1+1)$
Evolution Strategy takes a different approach: it maintains a single parent
and generates one child at each iteration. Its selection rule can be written
as
\begin{equation}\label{eq:standard_one_plus_one}
\theta^{(g+1)} =
\begin{cases}
\theta_{\text{child}}^{(g)} & \text{if } f(\theta_{\text{child}}^{(g)}) \leq f(\theta_{\text{parent}}^{(g)}), \\
\theta_{\text{parent}}^{(g)} & \text{otherwise}.
\end{cases}
\end{equation}

The $(1+1)$ structure is computationally simple because only one new
candidate has to be generated and compared with the current parent. It also
does not have the population-level averaging available in multi-parent
Evolution Strategies. Its replacement decision can therefore be sensitive
to errors in the fitness information \cite{Goldberg1989,Eiben2015}.

This property becomes particularly relevant when the individuals in the
evolutionary process are not numerical candidate solutions but complete
optimization algorithms. LLaMEA \cite{LLaMEA2024} applies this type of
program-space search to continuous black-box optimization. It embeds a
foundation model in an evolutionary loop in which generated optimization
programs are executed, evaluated, and iteratively modified.

Candidate algorithms in LLaMEA are represented as self-contained Python
classes following an object-oriented interface. Each candidate provides an
initialization method, \texttt{\_\_init\_\_}, which initializes parameters
and internal state, and an execution method,
\texttt{\_\_call\_\_(problem,budget)}, which receives a continuous black-box
objective and an evaluation budget and returns the best solution found by the
algorithm.

LLaMEA therefore turns algorithm design into a search over executable
programs. At generation $g$, the LLM acts as a semantic variation operator
and produces a child algorithm using information about the current parent,
previously evaluated candidates, and the optimization task.

For the discussion in this thesis, the generation context can be represented
compactly as
\begin{equation}\label{eq:llamea_context_assembly}
\mathcal{C}_g =
\left\{
\mathcal{P}_{\text{task}},
\mathcal{H}_{1 \dots g-1},
\mathcal{T}_{\text{incumbent}},
F_{\text{directive}}
\right\},
\end{equation}
where $\mathcal{P}_{\text{task}}$ specifies the optimization problem, search
bounds, evaluation budget, and required algorithm interface;
$\mathcal{H}_{1\dots g-1}$ stores information about previously generated
algorithms and their scores; $\mathcal{T}_{\text{incumbent}}$ contains the
current parent implementation and associated feedback; and
$F_{\text{directive}}$ specifies the requested modification.
This notation is a compact representation used in this thesis rather than a
formal notation introduced by the original LLaMEA paper.

The task description contains the continuous search domain, commonly
$[-5,5]^D$, together with the evaluation budget and interface requirements.
The history provides information about previously generated candidates and
their observed performance. The incumbent context provides the LLM with the
current algorithm and its evaluation feedback. The mutation directive can
ask the model to refine the existing strategy or to redesign it.

The generated response must be converted into executable source code before
evaluation. Code extraction removes conversational text surrounding the
generated program. The resulting candidate is then loaded into the execution
environment and evaluated using the optimization benchmark infrastructure.
In the standard LLaMEA setting, candidate algorithms are evaluated on
deterministic functions using \texttt{IOHexperimenter}
\cite{deNobel2024}. The evaluation tracks the objective gap
\begin{equation}
\Delta f(x_k)=f(x_k)-f(x^*)
\end{equation}
throughout the available evaluation budget and uses the BBOB target
$\Delta f \leq 10^{-8}$ as a high-precision reference
\cite{Hansen2010,LLaMEA2024}.

The original LLaMEA framework maps convergence behavior to a scalar fitness
measure so that candidates can be compared within the $(1+1)$ evolutionary
search \cite{LLaMEA2024}. Its normalized Area Over the Convergence Curve
(AOCC) can be written as
\begin{equation}\label{eq:aocc_metric_bg}
\text{AOCC}(y_{a,f}) =
\frac{1}{B}
\sum_{i=1}^{B}
\left(
1-
\frac{
\min(\max(y_i,\text{lb}),\text{ub})-\text{lb}
}{
\text{ub}-\text{lb}
}
\right),
\end{equation}
where
\begin{equation}
y_i =
\log_{10}
\left(
f(x_i)-f(x^*)
\right)
\end{equation}
represents the log-scaled distance to the optimum and the bounds are defined
by the evaluation range. In the original LLaMEA formulation,
the convergence values are log-scaled before AOCC is calculated, and the
resulting AOCC is aggregated over the BBOB functions and instances.
The mean AOCC across benchmark instances can be used as the scalar fitness for
candidate replacement in the original framework \cite{LLaMEA2024}.

Selection then follows the $(1+1)$ replacement rule:
\begin{equation}\label{eq:one_plus_one_selection}
\mathcal{A}_{\text{parent}}^{(g+1)} =
\begin{cases}
\mathcal{A}_{\text{child}}^{(g)} & \text{if } f(\mathcal{A}_{\text{child}}^{(g)}) \geq f(\mathcal{A}_{\text{parent}}^{(g)}), \\
\mathcal{A}_{\text{parent}}^{(g)} & \text{otherwise}.
\end{cases}
\end{equation}

The direction of this comparison depends on how the scalar fitness is
defined. For the original LLaMEA AOCC formulation, higher
AOCC indicates better anytime performance, so the child is retained when its
AOCC-based fitness is at least as large as that of the parent
\cite{LLaMEA2024}. The important point here is that one candidate is retained
as the parent, and this candidate becomes the starting point for the next
generation.

LLaMEA also incorporates feedback from the evaluation process. When a
generated algorithm is executed, its performance is recorded together with
the standard deviation obtained across repeated runs and any runtime error
information. For the elitist $(1+1)$ variant, this information is included in
the context used to generate the next candidate, together with the current
best-so-far algorithm and the history of previously generated algorithms
\cite{LLaMEA2024}. Runtime errors are handled explicitly: if an
evaluation fails, the candidate is assigned the worst fitness value in the
original maximization formulation. The feedback therefore allows the LLM to
use both performance information and execution errors when modifying or
redesigning an algorithm.

This feedback mechanism is relevant when considering LLaMEA under stochastic
fitness. The original framework was evaluated on noiseless BBOB functions,
so the fitness signal used for selection was not subject to the stochastic
evaluation considered in this thesis \cite{LLaMEA2024}. If the same elitist
$(1+1)$ selection mechanism is applied to noisy evaluations, an inferior
child can occasionally obtain a more favorable observed fitness than its
parent. Such a child may then be accepted as the new incumbent even though
its underlying performance is worse.

This creates a potential failure mode in the evolutionary search. Let
$\mathcal{A}_p$ denote the current parent algorithm and
$\mathcal{A}_c$ the generated child. Their observed performances can be
written as
\begin{equation}
\widehat{F}(\mathcal{A}) =
F(\mathcal{A})+\epsilon_{\mathcal{A}},
\end{equation}
where $F(\mathcal{A})$ denotes the underlying performance and
$\epsilon_{\mathcal{A}}$ represents the error introduced by stochastic
evaluation.

Under a favorable noise realization, a child with lower underlying
performance can therefore appear better than its parent. Strict elitist
selection can then replace the parent with this statistically favored
candidate. Since the selected candidate becomes the incumbent supplied to
the LLM for subsequent generations, the effect of the incorrect selection
can extend beyond the evaluation in which the ranking error occurred.

In a program-space search such as LLaMEA, this can influence which algorithmic
structure is used as the basis for subsequent modifications. If later
evaluations do not reproduce the same favorable noise realization, the
selected algorithm may no longer appear as strong as it did when it was
accepted. The search has nevertheless already moved to a different
incumbent. We refer to this potential failure mode as the
\emph{Ghost Parent} effect. This term describes the mechanism considered in
this thesis and is not a failure mode established by the original LLaMEA
study.

The distinction between algorithm evaluation and algorithm synthesis is
important here. Evaluating an already generated algorithm under noise
measures how its performance changes when the fitness signal is stochastic.
Running the synthesis process itself under noisy feedback introduces an
additional source of uncertainty: the stochastic fitness signal can influence
which generated algorithms are retained and subsequently used to guide
further LLM generations. The former therefore concerns the robustness of a
fixed algorithm, whereas the latter also concerns the reliability of the
algorithm-generation process under stochastic selection. This thesis
investigates the latter setting.

%------------------------------------------------------------------------------
\section{Research Gap}\label{sec:research_gap}
%------------------------------------------------------------------------------

The literature reviewed in this chapter covers two areas that are directly
relevant to this thesis. The first is continuous black-box optimization under
uncertainty. Previous studies have shown that noisy fitness can distort
candidate rankings, interfere with search directions, and affect adaptation
mechanisms in established optimization algorithms
\cite{Jin2005,Arnold2003,Hansen2009Noisy,Sarafian2024}. Noise-handling
strategies such as repeated evaluation and sample averaging can reduce the
effect of noisy observations, but they also require additional function
evaluations. In addition, they generally require the algorithm designer to
specify how uncertainty should be handled.

The second area is Automated Algorithm Design and the use of LLMs for program
synthesis. Earlier AAD approaches relied on predefined components or program
representations \cite{Burke2013,vanRijn2016,deNobel2021}. More recent
approaches use LLMs to generate and modify executable programs
\cite{Romera2023,Liu2023,Liu2024,ReEvo2024,LLaMEA2024}. LLaMEA is especially
relevant because it uses an LLM as a semantic variation operator inside a
$(1+1)$ evolutionary search and uses measured algorithm performance and
execution feedback to guide later generations \cite{LLaMEA2024}.

These two areas have not been studied to the same extent in combination. The
LLM-based AAD frameworks reviewed in this chapter primarily report evaluations
under deterministic or otherwise controlled benchmark conditions. The
comparison in Table~\ref{tab:llm_aad_comparison} shows that existing
approaches use different forms of performance, structural, reflective, or
problem-specific information, but stochastic fitness is not explicitly
treated as a source of uncertainty in the evolutionary selection process.

This creates a different problem from conventional noisy optimization. For an
established optimizer, noisy fitness affects the search over numerical
solutions. In an LLM-based AAD system, the same noisy feedback can also affect
the search over algorithmic programs. An incorrect selection can change the
parent program and, as a result, influence all subsequent generations.

The Ghost Parent mechanism described above illustrates one possible failure
mode. A favorable noise realization can cause an inferior generated algorithm
to become the incumbent. Subsequent generations are then conditioned on this
algorithm. The resulting behavior cannot be inferred from the deterministic
performance of the generated heuristics alone.

It therefore remains insufficiently characterized whether LLM-generated
continuous optimization heuristics remain effective under stochastic fitness
evaluation, how their performance changes as the noise level increases, and
how their behavior compares with established optimization methods under the
same conditions. It is also not yet clear to what extent the feedback-driven
generation process can produce algorithmic structures that remain effective
under noisy evaluations.

This thesis addresses this gap by studying LLM-based automated algorithm
design under controlled stochastic fitness evaluation. The investigation uses
continuous BBOB functions and compares LLM-generated heuristics with
established continuous optimization methods across different problem
dimensions and controlled noise levels. The following chapters describe the
stochastic evaluation framework, the LLaMEA-based synthesis procedure, and
the experimental protocol used to assess convergence, precision, and
robustness.
